% Draft: report subsection for the resource perspective (Sections 1.6-1.8).
% Bullet-point draft to be turned into prose and pasted into the ShareLaTeX
% report. Every number is reproducible from notebooks 01/02 and the
% models/*.json they emit. VIS-PLACEHOLDER bullets mark figures to export.

\subsection{Resource Availability, Permissions, and Allocation}

\textit{Johannes Terhorst}

The resource perspective adds three components on top of the engine, each
following the same \emph{fit offline, load at runtime} split as the rest of
the simulator:
\begin{itemize}
  \item Models are mined in notebooks (\texttt{01\_resource\_availability},
        \texttt{02\_resource\_permissions}) with pandas / scikit-learn /
        \texttt{ordinor} and serialised to small JSON parameter files in
        \texttt{models/}; the simulation only \emph{loads} JSON, so it stays
        dependency-light and the analysis stays reproducible.
  \item Each component sits behind a narrow interface
        (\texttt{PermissionModel}, an availability calendar object,
        \texttt{AllocationPolicy}), so a model can be swapped without touching
        the engine, the process component, or the event types.
\end{itemize}

\textbf{Resource availability (Section 1.6).}
\begin{itemize}
  \item \emph{Three log corrections first, each justified from the data:}
        (i) timestamps are UTC but working hours are local --- convert to
        \texttt{Europe/Amsterdam}, collapsing an apparent 1.17\,h summer/winter
        opening swing to 0.15\,h; (ii) only Workflow (\texttt{W\_}) events are
        human work with a real start/complete lifecycle (128{,}227
        \texttt{start} events), whereas \texttt{A\_}/\texttt{O\_} events only
        ever carry \texttt{complete} --- a naive ``filter to \texttt{complete}''
        would delete the entire signal; (iii) one ``resource'' is automated
        (\texttt{User\_1}, active on 99.2\% of calendar days, above the
        structural 71.4\%-weekday ceiling a human cannot exceed) and is
        excluded by a recomputed rule, not by name.
  \item \emph{Presence from the full lifecycle, bucketed by calendar date.} A
        resource is marked present on a day from its first to its last
        \texttt{W\_} signal \emph{that day}, using \emph{every} transition
        (schedule / start / suspend / resume / complete), not just start and
        complete. Each transition is counted on the date it actually occurred
        and never interpolated across a start-to-complete gap --- so a task
        started Friday and resumed Monday contributes a Friday point and a
        Monday point, and the intervening weekend gets nothing. Verified on the
        log: of 12{,}475 work items spanning $\geq$2 calendar days, across
        87{,}755 intervening days, \emph{zero} produced presence on a day with
        no real event. This is what keeps overnight and weekend-spanning tasks
        from fabricating availability.
  \item \emph{Limitation, tested --- automated lifecycle events.} The flip side
        of using the full lifecycle is that an \emph{automated} transition would
        still count as a presence signal: if the source system auto-suspended all
        open work at a fixed hour, that timestamp would enter the day's window
        even though the resource had stopped earlier, and the log could not tell
        a person closing a task from a nightly batch. We checked BPIC-17 for this
        and it does \emph{not} occur --- three fingerprints of an automated
        suspend are all absent: (i) the \texttt{suspend} hour-of-day curve tracks
        the \texttt{start} curve (peaks 9--11h, 14--16h; $\sim$0\% at 22:00),
        rather than spiking at one hour; (ii) timestamps are irregular
        (only 1.7\% on second~\texttt{:00}), not the round times a batch job
        produces; (iii) at most 2 cases ever share one exact \texttt{suspend}
        instant --- a ``suspend everything'' job would hit hundreds at once.
        So here \texttt{suspend} is a genuine human interaction and the workday
        is not inflated. The caveat is therefore general rather than active: on a
        different log with automated suspension this check would need repeating,
        and the $q_{90}$ window trim bounds the damage in any case.
  \item \emph{Two regimes in the lifecycle --- work vs waiting.} The intervals
        between \texttt{W\_} transitions are bimodal: active work
        (\texttt{start}/\texttt{resume}$\to$next) averages 5.8\,min and almost
        never spans a night (0.3\% cross midnight), whereas waiting
        (\texttt{suspend}$\to$\texttt{resume}, i.e.\ customer response) averages
        69.9\,h and crosses nights and weekends 55\% of the time. Across the log
        that is 24{,}634 active-hours against 15.05\,M waiting-hours ---
        \textbf{99.8\% of the elapsed \texttt{start}$\to$\texttt{complete} time
        is waiting, not work}. This is exactly why presence is reconstructed from
        dated transition \emph{points} rather than \texttt{start}$\to$%
        \texttt{complete} spans (which would drag weekend/overnight waiting into
        the workday); it is also a caveat for any processing-time model fitted on
        the raw span (cf.\ \S1.3), where the true active work is only minutes.
  \item \emph{\texttt{schedule} is a human signal, and harmless to availability.}
        \texttt{schedule} (a work item offered to a worklist) could be a system
        event, but it carries a resource 100\% of the time and in 85\% of cases
        the \emph{same} resource \texttt{start}s the item within one second --- a
        co-emitted marker, not an independent action (the other 13\% are a real
        worklist queue, p90 12\,h). Including it adds \emph{zero} new
        (resource, day) presence cells (+0.00\%), so the window model is robust
        to it either way.
  \item \emph{Basic --- weekly shift profile, in two tiers.} \emph{Tier A}
        (well-observed resources, $\geq$20 active days): a per-(resource,
        weekday) window fitted as $[q_{10}\text{ first event},\,
        q_{90}\text{ last event}]$ plus $P(\text{works that day})$. There is no
        fortnightly signal (odd/even-week split indistinguishable from a random
        split, $0.36$\,h vs $0.39$\,h median shift), so a \emph{weekly} cycle is
        tiled across the assignment's two-week interval rather than fitting
        noise into an A/B week.
  \item \emph{Tier B --- a coarser fit for sparse resources.} 27 humans have
        fewer than 20 active days --- too few distinct days to trust a shift
        \emph{per weekday}, but usually plenty of events (e.g.\ \texttt{User\_132}
        logs 2{,}622 events across 12 days). Rather than drop them --- which at
        runtime meant treating them as on shift \emph{24/7}, the opposite of
        sparse --- we fit a single \emph{pooled} window per resource from all
        its events and open it only on the weekdays it was actually seen. Weaker
        assumptions, but data-driven, and strictly better than an always-on
        fallback. After this, every human with any \texttt{W\_} signal has a
        window (148 of 148); only the genuine system account
        (\texttt{User\_1}, no \texttt{W\_} lifecycle) stays windowless and is
        modelled as always-on --- a batch process keeps no office hours. The
        \S1.6 model records which accounts are automated so the resource
        component can tell ``automated'' from ``unfitted'' rather than treating
        every windowless account as available around the clock.
  \item \emph{Why quantiles over min/max:} min/max reaches 100\% coverage but
        wastes half the modelled day (49.4\% utilisation) on outlier-stretched
        windows; $q_{10}/q_{90}$ covers \textbf{97.4\%} of real work events at
        57.2\% utilisation, against only \textbf{80.5\%} for a single global
        9--17 calendar --- the empirical case for per-resource shifts.
  \item \emph{Advanced --- year plan (\texttt{YearlyAvailability}).} The weekly
        profile plus (a) a discovered public-holiday calendar (weekdays of
        collective absence against a \emph{local} rolling baseline, recovering
        the Dutch bank holidays unprompted --- King's Day, Ascension, Whit
        Monday, Boxing Day --- and correctly omitting Good Friday); and (b)
        per-resource vacation, fitted only where identifiable.
  \item \emph{The vacation subtlety:} a flat ``$\geq 5$ absent days $=$ leave''
        rule yields an absurd 50 leave-days/resource/year, because the median
        resource works only $\sim$27\% of weekdays, so long absence runs occur
        $\sim$15$\times$/year by chance. Each resource gets its own threshold
        from its working rate; leave is fitted only for near-full-time staff
        ($\geq$60\%), where blocks then look like real annual leave (median
        10\,d, max 21\,d, $\sim$24\,days/year) --- a plausibility ceiling that
        \emph{emerges} from the data rather than being imposed.
  \item Holidays are a property of the \emph{calendar} (placed directly);
        vacation is a property of the \emph{person} (re-sampled per resource
        from its fitted distribution), so runs are statistically similar
        without replaying 2016's exact roster. Serialises to 81\,KB of
        parameters.
  \item \emph{Validation (second pass):} the simulated workforce reproduces the
        weekly rhythm, holiday troughs, and day-to-day variability (sd $6.8$
        vs real $7.2$). Folding in the sparse resources as Tier B narrows the
        mean-headcount gap from 9.2\% to \textbf{7.4\%} (model 37.3 vs real 40.2
        staff/weekday); the residual is a conservative error --- the model is
        slightly \emph{under}-staffed, so queues and waiting times are if
        anything overstated, not optimistic.
  \item \emph{VIS-PLACEHOLDER:} heat-map of the fitted weekly windows
        (open/close hour per resource $\times$ weekday) with a global 9--17
        band overlaid, showing how individual shifts diverge from office hours.
  \item \emph{VIS-PLACEHOLDER:} coverage vs.\ window-utilisation sweep
        (min/max, $q_{10}/q_{90}$, $q_{25}/q_{75}$, global 9--17) --- the table
        behind the window-fitting decision.
  \item \emph{VIS-PLACEHOLDER:} simulated vs.\ real weekday headcount over the
        log span, showing the reproduced holiday troughs and the deliberately
        un-modelled onboarding/truncation ramps at the edges.
\end{itemize}

\textbf{Resource permissions (Section 1.7).}

\emph{OrdinoR in brief (for readers new to it).} The basic model asks a
one-dimensional question --- ``has resource $r$ ever performed activity $a$?''
OrdinoR~\cite{yang2022ordinor} asks a richer one, built on three ideas:
\begin{itemize}
  \item \emph{Execution context.} Instead of the bare activity, each event is
        tagged with a triple (\emph{case type}, \emph{activity type},
        \emph{time type}) --- e.g.\ ``validate a \emph{car} loan on a
        \emph{Monday}''. Any dimension may be the wildcard $\bot$ (``any''), so
        the context can be as coarse as the activity alone or as fine as
        activity-\emph{and}-loan-type-\emph{and}-weekday.
  \item \emph{Groups and capabilities.} Resources that behave similarly are
        clustered into \emph{groups}; a group's \emph{capabilities} are the set
        of execution contexts its members work in; and a resource inherits the
        capabilities of every group it belongs to. Permission then becomes
        ``does $r$ belong to a group whose capabilities cover this context?''
        --- a \emph{role}, not a personal history. This is what lets the model
        grant a sensible pair a resource never personally did (its colleagues
        do, so the group does) and drop rare one-offs (a single stray event
        does not make it into the group's profile).
  \item \emph{Two quality numbers.} \emph{Fitness} = the share of real events
        the model allows (completeness); \emph{precision} = how tightly the
        allowed behaviour matches the log, i.e.\ how little extra it permits
        (exactness); \emph{F1} combines them. A model that permits everyone
        scores fitness 1 but near-zero precision --- a useless ``flower model''.
\end{itemize}
The three knobs we sweep map directly onto these ideas: how to \emph{define}
contexts (activity-only vs.\ adding case + time), how to \emph{cluster}
resources into groups (AHC $\to$ disjoint groups, MOC $\to$ overlapping), and
how to \emph{profile} a group (FullRecall = allow anything any member did $\to$
the flower model; OverallScore = keep only contexts that are frequent and
shared across the group).

\begin{itemize}
  \item \emph{Basic --- observed matrix.} A boolean
        resource\,$\times$\,activity matrix: permitted iff observed (the literal
        requirement). It is 63\% dense and pure \emph{memorisation} --- it
        cannot grant an unseen-but-sensible pair, a third of it rests on
        $\leq 10$ observations, and on a temporal hold-out it forbids
        \textbf{2{,}967 events that actually happen}.
  \item \emph{Advanced --- OrdinoR organizational model} (deployed model:
        \texttt{CT+AT+TT} contexts, AHC + complete linkage, OverallScore; 10
        groups over 144 resources). It both \emph{generalises} (grants sensible
        unseen pairs) and \emph{restricts} (drops rare one-offs), per the
        mechanism above.
  \item \emph{Design chosen by a 12-model sweep} ($3\times2\times2$: context
        definition $\times$ AHC/MOC discovery $\times$ FullRecall/OverallScore
        profiling): contexts from \texttt{CT+AT+TT}, \textbf{AHC} disjoint
        groups, \textbf{OverallScore} profiling (FullRecall is a ``flower
        model'' --- fitness 1.0, precision $\sim$0.01), and \textbf{complete}
        linkage (\texttt{ordinor} defaults to \texttt{single}, which chains
        102/144 resources into one blob). Reproduces the paper's BPIC-17
        experiment at F1 $\geq$ theirs on all AHC+OverallScore rows (best
        \textbf{0.744 vs 0.724}), with our fast conformance measures verified
        identical to \texttt{ordinor}'s to $10^{-9}$.
  \item \emph{Key design decision --- deploy the causally available model, not
        the best-scoring one.} The top-F1 model uses \emph{trace-clustered}
        case types (a summary of how a case behaved end-to-end), which cannot
        exist when a case needs its \emph{first} resource. We deploy the best
        model whose case types are known at arrival --- from the loan goal ---
        at a cost of 0.025 F1. (Conformance scores a model's \emph{description}
        of a finished log; a simulation needs a \emph{generator}.)
  \item \emph{Payoff in the simulation (the decisive test):} run under each
        model, the org model reproduces the real distribution of work across
        (resource, activity) pairs \textbf{best} (total-variation distance
        \textbf{0.575}, vs 0.696 for the observed matrix and 0.868 for the
        hardcoded top-20 baseline) --- and does so while being \emph{more}
        selective (about half as many resources permitted per activity). It
        wins by being more right, not more permissive.
  \item How the triple is enforced at runtime: the case type rides on the event
        payload (a loan goal sampled at arrival) and the time type comes from
        the clock the 1.6 calendar already needs, so both dimensions are
        available for free; a model with no opinion on a dimension degrades
        gracefully to a plain activity lookup.
  \item \emph{VIS-PLACEHOLDER:} fitness--precision scatter of the 12 swept
        models, highlighting the flower-model corner (FullRecall), the chosen
        \texttt{CT+AT+TT / AHC / OverallScore} point, and the deployed
        causally-available variant.
  \item \emph{VIS-PLACEHOLDER:} grouped bar of real-vs-simulated
        work-distribution TV-distance for the three permission models
        (hardcoded / observed / org), annotated with permitted
        resources-per-activity.
\end{itemize}

\textbf{Resource allocation and contention (Section 1.8 / Part~II baseline).}
\begin{itemize}
  \item \emph{Contention rework --- two-phase protocol.} Work items are
        \emph{requested}, then \emph{started}: the process component emits
        \texttt{ACTIVITY\_REQUEST} (``this item is enabled''); the resource
        component is the \emph{only} thing that emits \texttt{ACTIVITY\_START},
        and only once a resource is bound. The engine dispatches every event to
        all handlers unconditionally, so keeping a queued item as a request
        (invisible to the process component) is what stops it from being
        executed while it waits --- the fix for the case-forking bug --- and
        makes waiting time fall out as $\text{start}-\text{request}$.
  \item \emph{Random allocation (R-RBA, Russell et al.\ 2005 pattern~2).}
        Candidates are filtered by the permission model (\emph{who may}), then
        by live capacity (\texttt{capacity}$>1$ generalises the paper's
        single-slot resource) and on-shift status; among the survivors an
        \texttt{AllocationPolicy} picks one. The default \texttt{RandomPolicy}
        (R-RMA, pattern~15) is the assignment-mandated uniform random pick, and
        is the Part~II baseline every optimisation policy plugs into via the
        same seam without editing \texttt{ResourceComponent}.
  \item \emph{Saturation path (R-DE, pattern~19).} When every qualified resource
        is busy, the item is deferred onto a FIFO queue and re-offered the
        instant a qualified resource frees up (O(1) per release), rather than
        rejected.
  \item \emph{Adapting the heuristic to uncertain availability:} allocation
        reads the live \texttt{busy} counter, never a forecast, because
        stochastic service times make future availability unknown; the shift
        calendar gates \emph{allocation}, not \emph{execution}, so a task
        started at 16:55 finishes rather than being preempted at 17:00; and an
        activity nobody may perform runs unassigned and is counted in
        \texttt{stats()} rather than stranding the case forever.
  \item All picks use a seeded RNG, so the whole allocation path is bit-for-bit
        reproducible under a fixed seed, as the grading requires.
  \item \emph{VIS-PLACEHOLDER:} sequence diagram of the two-phase protocol ---
        \texttt{ACTIVITY\_REQUEST} $\rightarrow$ permission/capacity/shift
        filter $\rightarrow$ policy pick $\rightarrow$ \texttt{ACTIVITY\_START},
        with the R-DE deferral loop back through \texttt{RESOURCE\_AVAILABLE}.
\end{itemize}
